\documentclass[10pt]{article}
\usepackage[a4paper,left=2.5cm,right=2.5cm,top=2cm,bottom=2.5cm]{geometry}
\usepackage{graphicx}
\graphicspath{{imagenes/}}
\usepackage[LGR,T1]{fontenc}
\usepackage[greek,spanish]{babel}
\usepackage{mwe}
\usepackage{subfiles}
\usepackage{noto}
\usepackage{parskip}
\usepackage[backend=biber, style=authoryear]{biblatex}
\addbibresource{biblio.bib}
\setcounter{biburlnumpenalty}{100}
\setcounter{biburlucpenalty}{100}
\setcounter{biburllcpenalty}{100}
\usepackage{pgfgantt}
\usepackage[most]{tcolorbox}

% comentarios de revisión, visibles en el PDF como recuadros de color;
% cambiar a \mostrarcomentariosfalse para ocultarlos en la versión final
\newif\ifmostrarcomentarios
\mostrarcomentariostrue
\ifmostrarcomentarios
	\NewDocumentEnvironment{comentario}{m}{%
		\begin{tcolorbox}[breakable, colback=yellow!12, colframe=orange!70!black,
			fonttitle=\bfseries\small, fontupper=\small\itshape,
			title={Comentario de #1}]%
	}{\end{tcolorbox}}
\else
	\NewDocumentEnvironment{comentario}{m+b}{}{}
\fi
% marcadores (outline) del PDF; hyperref va al final del preámbulo
\usepackage[hidelinks,bookmarksopen,bookmarksnumbered,pdfpagemode=UseOutlines]{hyperref}

% marcador para capítulos en borrador
\newcommand{\esborrador}{}

% aviso al inicio de los capítulos que son esbozos en escritura
\newcommand{\avisoesbozo}{%
	\begin{tcolorbox}[colback=gray!8, colframe=gray!60!black,
		fonttitle=\bfseries\small, fontupper=\small,
		title={Aviso}]
		Este capítulo es un esbozo inicial de la tesis, actualmente en escritura. Su contenido es preliminar e incompleto, y se incluye como muestra del trabajo en curso.
	\end{tcolorbox}%
}

% \capitulo{nombre} incluye capitulos/nombre.tex,
% excepto si el archivo contiene \esborrador
\ExplSyntaxOn
\NewDocumentCommand{\capitulo}{m}{
	\file_get:nnN {capitulos/#1.tex} {} \l_tmpa_tl
	\tl_if_in:NnF \l_tmpa_tl {\esborrador} { \subfile{capitulos/#1.tex} }
}
\ExplSyntaxOff

\begin{document}

\begin{titlepage}
	\centering
	{\LARGE \textsc{Universidad de Santiago de Chile}\par}
	\vspace{1cm}
	{\Large \textsc{Instituto de Estudios Avanzados}\par}
	\vspace{1.5cm}
	{\huge\bfseries Formulario de proyecto de tesis\par}
	\vspace{2cm}
	{\huge\bfseries Popusintesíntesis\par}
	\vspace{2cm}
	{\Large Aarón Montoya Moraga\par}
	\vfill
	supervisión por \par
	Dra. Cynthia Shuffer

	\vfill

% Bottom of the page
	{\large \today\par}
\end{titlepage}

\newpage

\tableofcontents

\newpage

\subfile{secciones-formulario/resumen}
\newpage
\subfile{secciones-formulario/palabras-clave}
\newpage
\subfile{secciones-formulario/hipotesis}
\newpage
\subfile{secciones-formulario/objetivos}
\newpage
\subfile{secciones-formulario/investigacion-propuesta}
\newpage
\subfile{secciones-formulario/metodologia}
\newpage
\subfile{secciones-formulario/formatos}
\newpage
\subfile{secciones-formulario/trabajo-avanzado}
\newpage
\subfile{secciones-formulario/plan-trabajo}
\newpage
\subfile{secciones-formulario/indice-tentativo}
\newpage

\capitulo{01-sintesis-sonora}
\capitulo{02-sintesintesis-sonora}
\capitulo{03-popusintesintesis-sonora}
\capitulo{04-modulos}
\capitulo{05-metaforas}
\capitulo{06-taxonomias}
\capitulo{07-relojes}
\capitulo{08-secuenciadores}
\capitulo{09-osciladores}
\capitulo{10-filtros}
\capitulo{11-envolventes}
\capitulo{12-amplificadores}
\capitulo{13-moduladores}
\capitulo{14-efectos}
\capitulo{15-parlantes}
\capitulo{16-controladores}
\capitulo{17-interfaces}
\capitulo{18-soldaduras}
\capitulo{19-variables}
\capitulo{20-funciones}
\capitulo{21-clases}
\capitulo{22-bibliotecas}
\capitulo{23-documentacion}
\capitulo{24-prototipado}
\capitulo{25-conectores}
\capitulo{26-carcasas}
\capitulo{27-manuales}
\capitulo{28-reparaciones}
\capitulo{29-mantenimientos}
\capitulo{30-partituras}

\subfile{secciones-formulario/fuentes}

\newpage

\subfile{secciones-formulario/apendice}

\end{document}